% board: Nexys A7-100T I/O map -- slide switches, LEDs, 8-digit display, push buttons and the 100 MHz clock around the FPGA that holds the LS core
\begin{tikzpicture}[x=1mm,y=1mm,
  bd txt/.style={font=\sffamily\scriptsize,text=lsink,inner sep=0pt,anchor=base},
  bd note/.style={bd txt,text=lsgrey},
  bd brace/.style={decorate,decoration={brace,amplitude=1.2mm},draw=lsgrey,line width=0.5pt},
  bd sw/.style={draw=lsnavy!80,fill=white,line width=0.5pt,rounded corners=0.6pt},
  bd knob/.style={fill=lsnavy,rounded corners=0.4pt},
  bd led/.style={circle,inner sep=0pt,minimum size=2.9mm,line width=0.5pt},
  bd pill/.style={fill=lsamber,text=white,rounded corners=1.5pt,inner xsep=3pt,inner ysep=1.2pt,
    font=\sffamily\scriptsize\bfseries},
]
\def\bdX#1{81+(7.5-#1)*9.2}% centre x of LED #1 and switch #1
\pgfmathparse{Hex(\cycS)}\let\bdHex\pgfmathresult

% ---- board ----
\filldraw[fill=lspaper,draw=lsnavy,line width=0.5pt,rounded corners=3mm] (0.5,0.5) rectangle (161.5,91.5);
\foreach \bdx/\bdy in {3.8/3.8,158.2/3.8,3.8/88.2,158.2/88.2}
  \filldraw[fill=white,draw=lsrule,line width=0.5pt] (\bdx,\bdy) circle (1.2);
\node[bd note,font=\sffamily\scriptsize\bfseries,anchor=base west] at (7,85.2) {Nexys A7-100T};

\begin{scope}[yshift=1mm]% lower block lifted so the bottom inset matches the top
% ---- slide switches SW15..SW0, drawn as set for S64 with view 4 (only SW10 up) ----
\foreach \bdk in {0,...,15} {
  \filldraw[bd sw] ({\bdX{\bdk}-2.2},11.5) rectangle ({\bdX{\bdk}+2.2},19.9);
  \ifnum\bdk=10 \fill[bd knob] ({\bdX{\bdk}-1.5},15.9) rectangle ({\bdX{\bdk}+1.5},19.2);
  \else \fill[bd knob] ({\bdX{\bdk}-1.5},12.2) rectangle ({\bdX{\bdk}+1.5},15.5); \fi
  \node[bd note] at ({\bdX{\bdk}},21.8) {\bdk};
}
\foreach \bdl/\bdr in {15/13,12/11,10/8,7/0}
  \draw[bd brace,decoration={mirror}] ({\bdX{\bdl}-2.6},10.2) -- ({\bdX{\bdr}+2.6},10.2);
\node[bd txt]  at ({\bdX{14}},6.0)   {program select};
\node[bd note] at ({\bdX{14}},3.1)   {read at reset};
\node[bd txt]  at ({\bdX{11.5}},6.0) {run mode};
\node[bd txt]  at ({\bdX{9}},6.0)    {display view};
\node[bd txt]  at ({\bdX{3.5}},6.0)  {IN data};
\node[bd note] at ({\bdX{3.5}},3.1)  {SW7..6 = image};

% ---- LEDs LED15..LED0 ----
\foreach \bdk/\bdc/\bdf in {15/lsblue/lsblue!18,14/lsblue!65/lsblue!12,13/lsteal/lsteal!18,
    12/lsteal!65/lsteal!12,11/lsamber/lsamber!20,10/lsrose!80/lsrose!15}
  \node[bd led,draw=\bdc,fill=\bdf] at ({\bdX{\bdk}},26.8) {};
\foreach \bdk in {0,...,9}
  \node[bd led,draw=lsgrey,fill=white] at ({\bdX{\bdk}},26.8) {};
\foreach \bdl/\bdr in {15/12,11/10,9/8,7/0}
  \draw[bd brace] ({\bdX{\bdl}-2.2},29.5) -- ({\bdX{\bdr}+2.2},29.5);
\node[bd txt,text=lsblue]            at ({\bdX{13.5}},35.2) {FETCH DECODE};
\node[bd txt,text=lsteal]            at ({\bdX{13.5}},32.3) {EXECUTE WRITEBACK};
\node[bd txt,text=lsamber!90!black]  at ({\bdX{10.5}},35.2) {STREAM};
\node[bd txt,text=lsrose]            at ({\bdX{10.5}},32.3) {DRAIN};
\node[bd txt]                        at ({\bdX{8.5}},32.3)  {V1 V2};
\node[bd txt]                        at ({\bdX{3.5}},32.3)  {OUTR (the neuron's output byte)};
\end{scope}

% ---- 8-digit seven-segment display, showing view 4 (cycle counter) after S64 ----
\def\bdW{1.7}\def\bdH{3.0}\def\bdT{0.72}\def\bdG{0.2}
\def\bdHs#1#2{\fill[#2] ({-\bdW+\bdG},#1) -- ({-\bdW+\bdG+\bdT/2},{#1+\bdT/2})
  -- ({\bdW-\bdG-\bdT/2},{#1+\bdT/2}) -- ({\bdW-\bdG},#1) -- ({\bdW-\bdG-\bdT/2},{#1-\bdT/2})
  -- ({-\bdW+\bdG+\bdT/2},{#1-\bdT/2}) -- cycle;}
\def\bdVs#1#2#3{\fill[#3] (#1,{#2+\bdG}) -- ({#1+\bdT/2},{#2+\bdG+\bdT/2})
  -- ({#1+\bdT/2},{#2+\bdH-\bdG-\bdT/2}) -- (#1,{#2+\bdH-\bdG}) -- ({#1-\bdT/2},{#2+\bdH-\bdG-\bdT/2})
  -- ({#1-\bdT/2},{#2+\bdG+\bdT/2}) -- cycle;}
\def\bdSeg#1{\ifnum#1=1 \def\bdc{lsamber}\else\def\bdc{lsink!82!white}\fi}
\def\bdSegs#1#2#3#4#5#6#7{%
  \bdSeg{#1}\bdHs{\bdH}{\bdc}\bdSeg{#2}\bdVs{\bdW}{0}{\bdc}\bdSeg{#3}\bdVs{\bdW}{-\bdH}{\bdc}%
  \bdSeg{#4}\bdHs{-\bdH}{\bdc}\bdSeg{#5}\bdVs{-\bdW}{-\bdH}{\bdc}\bdSeg{#6}\bdVs{-\bdW}{0}{\bdc}%
  \bdSeg{#7}\bdHs{0}{\bdc}}
\def\bdDigit#1#2#3{% centre x, centre y, value 0..15; segments a..g
  \begin{scope}[shift={({#1},{#2})},xslant=0.1]
    \ifcase#3 \def\bdS{1111110}\or\def\bdS{0110000}\or\def\bdS{1101101}\or\def\bdS{1111001}%
    \or\def\bdS{0110011}\or\def\bdS{1011011}\or\def\bdS{1011111}\or\def\bdS{1110000}%
    \or\def\bdS{1111111}\or\def\bdS{1111011}\or\def\bdS{1110111}\or\def\bdS{0011111}%
    \or\def\bdS{1001110}\or\def\bdS{0111101}\or\def\bdS{1001111}\or\def\bdS{1000111}\fi
    \expandafter\bdSegs\bdS
    \fill[lsink!82!white] ({\bdW+0.95},{-\bdH}) circle (0.38);
  \end{scope}}
\foreach \bdm in {0,1}
  \fill[lsink,rounded corners=0.8mm] ({7+\bdm*26.7},68.8) rectangle ++(24.5,10.7);
\foreach \bdj in {0,...,7} {
  \pgfmathtruncatemacro\bdv{\bdj<5 ? 0 : (\bdj==5 ? int(\cycS/256) : (\bdj==6 ? mod(int(\cycS/16),16) : mod(\cycS,16)))}
  \bdDigit{10.6+\bdj*5.5+(\bdj>3)*4.7}{74.15}{\bdv}
}
\node[bd txt]  at (32.6,64.6) {8-digit display: view selected by SW10..8};
\node[bd note] at (32.6,61.7) {shown: view 4, S64 takes \pred{\cycS{} = 0x\bdHex} cycles};

% ---- 100 MHz oscillator into pin E3 ----
\draw[lsgrey,line width=0.5pt] (60.6,51.5) -- (66,51.5);
\filldraw[fill=white,draw=lsgrey,line width=0.5pt,rounded corners=0.8pt] (54.2,49.5) rectangle (60.6,53.5);
\draw[lsgrey,line width=0.4pt] (55,50.4) -- ++(0.8,0) |- ++(1.2,2.2) |- ++(1.2,-2.2) |- ++(1.2,2.2) |- ++(0.4,-2.2);
\node[bd txt,anchor=east] at (53.2,51.5) {100 MHz (E3)};

% ---- FPGA ----
\fill[lsnavy,rounded corners=1.2mm] (66,49.5) rectangle (96,79.5);
\draw[lsgrey!70,line width=0.4pt,rounded corners=0.8mm] (67.4,50.9) rectangle (94.6,78.1);
\fill[lsgrey] (69.4,76) circle (0.7);
\node[bd txt,text=lsrule] at (81,70) {Artix-7};
\node[bd txt,font=\sffamily\footnotesize\bfseries,text=white] at (81,65.3) {XC7A100T};
\node[bd pill] at (81,57.5) {LS core};

% ---- push buttons ----
\foreach \bdx/\bdy/\bdn/\bdf/\bdon in {128.75/78.5/BTNU/reset/1,114.75/64.5/BTNL/unused/0,
    128.75/64.5/BTNC/{go / step}/1,142.75/64.5/BTNR/unused/0,128.75/50.5/BTND/unused/0} {
  \ifnum\bdon=1 \def\bdc{lsnavy}\def\bdt{lsink}\else\def\bdc{lsrule}\def\bdt{lsgrey}\fi
  \filldraw[fill=white,draw=\bdc,line width=0.5pt,rounded corners=0.8pt] (\bdx-2.8,\bdy-2.8) rectangle (\bdx+2.8,\bdy+2.8);
  \fill[\bdc] (\bdx,\bdy) circle (1.9);
  \node[bd note] at (\bdx,\bdy-5.8) {\bdn};
  \node[bd txt,text=\bdt] at (\bdx,\bdy-8.7) {\bdf};
}
\end{tikzpicture}%
